\subsubsection{Microsoft Dynamics 365 for Finance and Supply Chain Management}
\ac{D365FnSCM}, früher auch Dynamics 365 for Operations genannt, ist ein cloudbasiertes \ac{ERP}-System von Microsoft, welches eine Vielzahl an Geschäftsprozessen eines Unternehmens in einem System bündelt. Es unterstützt dabei verschiedene Bereiche wie Finanzen, Vertriebs- und Einkaufsprozesse, Produktions- und Lagerverwaltung, Projektmanagement und vieles mehr. Als Teil der Dynamics 365-Produktlinie richtet sich diese Lösung insbesondere an mittelständische und große Unternehmen und integriert Finanz- und Operations-Funktionen nahtlos auf einer Plattform. \cite[\pagef 19ff]{mohta_implementing_2017}

%\subsubsection{Angebotsprozess}

%Einer dieser Geschäftsprozesse, die von \ac{D365FnSCM} unterstützt werden, ist der Angebotsprozess. dieser umfasst die Erstellung, Verwaltung und Nachverfolgung von Angeboten für bestehende oder potenzielle Kunden. In der ERP-Anwendung wird ein Verkaufsangebot für einen Interessenten oder Kunden erstellt, in dem Produkte, Mengen, Preise, Rabatte und Lieferdaten festgehalten werden. Dieses Angebotsdokument stellt einen formalen Vorschlag an den Kunden dar. Wird das Angebot vom Kunden angenommen und bestätigt, so wird es im System in einen Verkaufsauftrag umgewandelt. Der Fokus dieser Arbeit liegt dabei nur auf der Erstellung eines Verkaufsangebots, wobei sich der Erstellungsprozess eines Angebots aus mehreren Schritten zusammensetzt, wie in Abbildung \ref{fig:angebotserstellungsprozess} dargestellt. % Als Lesezeichen markiert

%Hier BPMN wie man ein Angebot erstellt in D365FO

\subsubsection{Microsoft Dataverse und Power Platform}

Eine zentrale Komponente für die Integration von Dynamics-365-Anwendungen ist Microsoft Dataverse. Dataverse ist eine cloudbasierte Datenplattform, die als gemeinsame Datenbasis für Dynamics-365-Anwendungen dient und die strukturierte Verwaltung unternehmensweiter Geschäftsdaten ermöglicht. Sämtliche Dynamics-365-Daten, einschließlich der Daten aus \ac{D365FnSCM}, können in Dataverse bereitgestellt werden. Durch das einheitliche Datenmodell unterstützt Dataverse die Integration und Erweiterung von Dynamics-365-Anwendungen sowie die Anbindung weiterer Dienste innerhalb des Microsoft-Ökosystems. Auf dieser Grundlage lassen sich unter anderem Anwendungen, automatisierte Abläufe und \ac{KI}-Agenten realisieren, die direkt auf die zentralen Unternehmensdaten zugreifen. \cite{milindav_abrufen_nodate}


Damit Daten aus \ac{D365FnSCM} in der Power-Platform-Umgebung genutzt und verfügbar gemacht werden können, muss die Power Platform Integration für die D365-Umgebung aktiviert sein \cite{cabeln_enable_nodate}. Diese Kopplung verbindet die \ac{ERP}-Datenbank mit einer Dataverse-Datenbank. In einer solchen verknüpften Umgebung können Änderungen in \ac{D365FnSCM} nahezu in Echtzeit mit Dataverse synchronisiert werden und umgekehrt. Damit dies funktioniert, stellt Microsoft dafür den sogenannten Dual-write-Dienst bereit. Dual-write ist eine vorkonfigurierte Integrationsfunktion, die eine bidirektionale, nahezu echtzeitfähige Synchronisation zwischen \ac{D365FnSCM} und Dataverse ermöglicht. Sobald Dual-write konfiguriert ist, werden Änderungen an Geschäftsdaten in \ac{D365FnSCM} automatisch an Dataverse weitergegeben und umgekehrt. Dies sorgt für konsistente Daten über alle angebundenen Systeme hinweg und bildet die Grundlage, um D365 Anwendungen mit den ERP-Daten zu versorgen. \cite{ramakrishnamoorthy_dual-write_nodate}

Alternativ zur physischen Datensynchronisation mittels Dual-write bietet \ac{D365FnSCM} auch die Möglichkeit, sogenannte virtuelle Entitäten in Dataverse bereitzustellen. Dabei handelt es sich um eine direkte Anbindung der ERP-Datenbank an Dataverse, ohne die Daten zu duplizieren. Viele Entitäten der D365-Apps sind über OData-Webservices verfügbar, jedoch nicht standardmäßig als Dataverse-Entitäten sichtbar. Nach Verknüpfung der Umgebung kann durch einen Administrator festgelegt werden, welche dieser ERP-Entitäten in Dataverse sichtbar gemacht werden sollen. Wird eine Entität in Dataverse veröffentlicht, entsteht dort eine entsprechende virtuelle Entität, welche eine Tabelle darstellt, die in Echtzeit auf die zugrunde liegenden Daten in \ac{D365FnSCM} zugreift. Auf diese Weise kann ein Datensatz aus \ac{D365FnSCM} als Dataverse-Tabelle gelesen und durch D365-Apps verwendet werden, ohne dass eine separate Kopie des Datensatzes gespeichert wird. \cite{jaredha_enable_nodate}

\subsubsection{Microsoft Copilot Studio}

Microsoft verfolgt das Ziel, generative \ac{KI} umfassend in Geschäftsprozesse zu integrieren. In \ac{D365FnSCM} selbst sind erste Copilot-Funktionen, wie beispielsweise ein Hilfefenster mit einem Chatbot oder automatisierte Zusammenfassungen an bestimmten Formularen eingebettet. Diese Funktionen greifen auf große \ac{LLM}s zurück, die über Azure gehostet und über Dataverse verbunden sind, und nutzen die im System vorhandenen Daten als Kontext. \cite{cabeln_enable_nodate-1}

Für verschiedenste unternehmensspezifische Anwendungsfälle reichen die angebundenen Copilot-Funktionen jedoch nicht aus. Microsoft Copilot Studio ist daher eine Plattform, mit der eigene \ac{KI}-gestützte Agenten für individuelle Szenarien erstellt und verwaltet werden können. Es handelt sich um ein Low-Code-Werkzeug, das die Konfiguration solcher Agenten ohne tiefe \ac{KI}- oder Programmierkenntnisse ermöglicht. Ein Copilot-Studio-Agent ist ein \ac{KI}-basierter digitaler Assistent, der in natürlicher Sprache mit Benutzern interagiert und definierte Aufgaben innerhalb von Geschäftsprozessen unterstützt. \cite{iaanw_overview_nodate}
Die konzeptionellen Grundlagen und Funktionsweisen von \ac{KI}-Agenten werden im weiteren Verlauf dieses Kapitels detailliert erläutert.

In Copilot Studio können verschiedene Komponenten konfiguriert werden, um den Agenten an die spezifischen Anforderungen eines Anwendungsfalls anzupassen. 

Die erste Komponente sind die sogenannten Instructions. Durch diese Anweisungen wird das Verhalten des Agenten definiert, indem Rollen, Richtlinien oder Gedächtnisinformationen festgelegt werden. Sie dienen als Leitfaden für den Agenten und beeinflussen, wie dieser auf Benutzereingaben reagiert. \cite{noauthor_copilot_nodate}

Eine weitere Komponente sind die Topics, welche einen thematischen Dialogbaustein eines Agenten darstellen. Ein Topic definiert dabei einen abgeschlossenen Gesprächsablauf zu einem bestimmten Anliegen oder Anwendungsfall. Diese Topics definieren dann innerhalb dieses Ablaufs wie eine Unterhaltung verlaufen soll. Diese Topics können durch sogenannte Trigger ausgelöst werden, also bestimmte Schlüsselwörter oder Phrasen, die der Benutzer eingibt und innerhalb dieser Topics sind Konversationsknoten enthalten, welche dann die eigentlichen Aktionen des Topics darstellen. Diese Knoten können verschiedene Arten von Knoten unterteilt werden. Es gibt beispielsweise Frage-Knoten, um Benutzereingaben abzufragen, Bedingungen für Verzweigungen oder Aktions-Knoten, um Tools aufzurufen, welche im späteren genauer erläutert werden. Insgesamt ist ein Topic also eine Sequenz von Knoten mit verschiedenen Aufgaben. \cite{henryjammes_topics_nodate}

Die nächste Komponente sind die Tools oder auch Actions genannt. Diese Tools ermöglichen es dem Agenten, über die reine Konversation hinaus in externen Systemen Aktionen auszuführen, um die Benutzeranfragen direkt zu bearbeiten. Für diese Aktionen können verschiedene Arten von Tools hinterlegt werden. Beispiele hierfür sind das Versenden von E-Mails, Erstellen von Kalendereinträgen, aber auch das Bearbeiten, Abrufen oder Erstellen von Datensätzen in Dataverse, wodurch dann auch das Lesen und Schreiben von Datensätzen in \ac{D365FnSCM} ermöglicht wird. Alle einem Agenten hinzugefügten Tools erscheinen anschließend in einer Tool-Liste, von wo aus sie entweder für den ganzen \ac{KI}-Agenten zur Verfügung gestellt werden können, wenn der generative Modus ausgewählt ist, bei dem der \ac{KI}-Agent selbst entscheidet, wann er das Tool nutzt, aber auch nur explizit innerhalb eines bestimmten Topics aufgerufen werden können. \cite{garypretty_add_nodate}

Eine wichtige grundlegende Komponente sind die Wissensquellen, die dem Agenten als Kontext dienen und ermöglichen es Unternehmensdaten und Dokumentationen mit dem Agenten zu verknüpfen. \cite{mssqlgirl_knowledge_nodate} Durch diese Daten ist es dann dem \ac{KI}-Agenten möglich, Antworten mit Bezug auf die spezifischen Unternehmensinformationen zu geben \cite{iaanw_overview_nodate}. Diese Daten können aus verschiedensten Quellen stammen. Beispiele hierfür sind SharePoint-Inhalte, öffentliche oder interne Websites oder auch Dynamics 365-Daten. Diese Wissensquellen lassen sich sowohl global für einen Agenten, als auch an bestimmten Topics hinterlegen. \cite{mssqlgirl_knowledge_nodate}

Bei diesen Wissensquellen ist es auch möglich, das Dataverse als Quelle zu verwenden. Das bedeutet, dass Daten aus \ac{D365FnSCM}, die in Dataverse als virtuelle Entitäten vorliegen, direkt vom Agenten abgefragt und genutzt werden können. \cite{relnotes_add_nodate}

